\documentclass[12pt,a4paper]{article}
        \makeatletter
        \def\input@path{{../}}
        \makeatother
        \usepackage{almeria-fichas}

        \FichaMeta{Sesion 05}{Rectas, segmentos y angulos}{Bloque 2 · Geometria urbana}{Identificar rectas, segmentos, paralelas, perpendiculares y angulos en un plano urbano.}{Conceptual}{Recordar, Comprender, Aplicar}

        \begin{document}

        \FichaCabecera

        \begin{materialesbox}
        \begin{itemize}
    \item Ficha impresa
    \item Lapiz
    \item Regla
    \item Lapis azul, rojo y verde
\end{itemize}
        \end{materialesbox}

        \begin{pasosbox}
        \begin{enumerate}
    \item Lee primero los dibujos o ejemplos.
    \item Marca con color antes de escribir la explicacion.
    \item Si aparece un angulo, indica tambien su tipo.
    \item No olvides nombrar los extremos de un segmento.
\end{enumerate}
        \end{pasosbox}

        \begin{recordatoriobox}
        \begin{itemize}
    \item Una recta no tiene extremos. Un segmento si.
    \item Las rectas paralelas no se cortan.
    \item Las rectas perpendiculares forman 90º.
\end{itemize}
        \end{recordatoriobox}

        \begin{apoyobox}
        \begin{itemize}
    \item Usa azul para paralelas, rojo para perpendiculares y verde para segmentos.
    \item Si ves una calle real, piensa si tiene inicio y final: entonces es un segmento.
    \item Puedes hacer un pequeño dibujo al lado de la respuesta.
\end{itemize}
        \end{apoyobox}

        \BloomSection{recordar}

\BloomExercise{recordar}{1}{Conceptos basicos}{Define brevemente: \textbf{recta}, \textbf{segmento}, \textbf{angulo recto}.}{5}

\BloomExercise{recordar}{2}{Clasifica}{Indica si estos ejemplos representan recta o segmento: borde de una mesa, tramo de calle, linea matematica, lado de un folio.}{5}

\BloomSection{comprender}

\BloomExercise{comprender}{1}{Paralelas o perpendiculares}{Explica la diferencia entre \textbf{paralelas} y \textbf{perpendiculares} usando una frase y un ejemplo de ciudad.}{6}

\BloomExercise{comprender}{2}{Tipo de angulo}{Explica por que un angulo de 120º es \textbf{obtuso} y uno de 55º es \textbf{agudo}.}{6}

\BloomSection{aplicar}

\BloomExercise{aplicar}{1}{Plano anotado}{Dibuja un pequeno plano con \textbf{dos calles paralelas}, \textbf{una interseccion perpendicular} y \textbf{un segmento} con extremos $A$ y $B$.}{8}

\BloomExercise{aplicar}{2}{Marca elementos}{En una esquina de barrio inventada, señala \textbf{dos angulos} distintos y escribe su tipo: agudo, recto u obtuso.}{7}

        \begin{evidenciabox}
        \begin{itemize}
    \item Distinguir recta, semirrecta y segmento.
    \item Reconocer angulos agudos, rectos y obtusos en la ciudad.
    \item Justificar si dos calles son paralelas o perpendiculares.
\end{itemize}
        \end{evidenciabox}

        \end{document}
